\printedpage{1}
% head: A. Grothendieck : Dérivateurs, Ch. XVIII, Édition des Universités de Montpellier II et Paris VII
% foot: Version 17 décembre 2025

\msp{page 50}

\section*{4\quad Paratopoi and the operation $\mathcal{M}\boxtimes\mathcal{N}$}

\subsection*{4.1\quad Pseudo-topoi. Problem of the product.}

I call \textit{pseudo-$\mathfrak{U}$-topos category}, or simply \textit{pseudo-topos}, a $\mathfrak{U}$-category $\mathcal{M}$ stable under $\mathfrak{U}$-direct limits, and which admits a \textit{small} subcategory $\mathcal{M}_0$ which is \textit{strictly generating} (i.e. generating by strict epimorphisms). Then $\mathcal{M}$ also admits $\mathfrak{U}$-inverse limits (SGA 4 I 8.12.7). (But the example of (Gr), cf. loc. cit. 8.12.9, shows that $\mathcal{M}^\circ$ is not necessarily, for all that, a pseudo-$\mathfrak{U}$-topos, for lack on the part of $\mathcal{M}$ of admitting a small strictly \textit{co}generating subcategory. But theorem 2.6 \texttt{[of chapter XIX]} gives conditions on $\mathcal{M}$ (mainly of exactness type) % [sic] par forcément -> pas forcément, misprint; also the parenthesis opened at "(Mais" is never closed
\msp{page 51}
for $\mathcal{M}$ to admit enough injectives -- which will then imply the existence of a small strictly cogenerating subcategory, hence that $\mathcal{M}^\circ$ is also a pseudo-topos (\footnote{This is the case in particular if $\mathcal{M}$ is a topos.}).

I visualize pseudo-topos categories as being the category of ``sheaves'' on an object $\mathcal{X}$, visualized as a kind of ``space'' which I distinguish from the category $\mathcal{M}=\Fais(\mathcal{X})$ of sheaves on $\mathcal{X}$. It is $\mathcal{X}$ which will be for this the ``pseudo-topos''. The pseudo-topoi form a 2-category, the variance in $\mathcal{X}$ being opposite to that in $\Fais(\mathcal{X})=\mathcal{M}$. Precisely, if $\mathcal{X}$ corresponds to $\mathcal{M}$, $\mathcal{Y}$ to $\mathcal{N}$, a morphism $f$ from $\mathcal{X}$ to $\mathcal{Y}$ will be given by a functor
\[
f^* \;:\; \mathcal{N} \longrightarrow \mathcal{M}
\tag*{(4.1.1)}
\]
in the opposite direction, commuting with arbitrary direct $\varinjlim$ (\footnote{Say here that the objects of $\mathcal{M}$ are visualized as generalized ``opens'' of $\mathcal{X}$. We have $\mathcal{M}\xrightarrow{\sim}\uHom^!(\mathcal{M}^\circ,\Ens)$, i.e. $\mathcal{M}$ is indeed interpreted as the ``sheaves'' on this category $\mathcal{M}$ of opens \ldots}). (\textbf{N.B.} if $\mathcal{X}$, $\mathcal{Y}$ are even topoi, i.e. $\mathcal{M}$, $\mathcal{N}$ topos-categories, and if we want $f$ to be a morphism of topoi, we will require moreover that $f^*$ be moreover \textit{left exact}. We will say that a \textit{morphism of pseudo-topoi is exact} if moreover $f^*$ is left exact.)

\msp{page 52}
By virtue of the dual of loc. cit. 8.12.7, this is equivalent to the fact that $f^*$ admits a right adjoint $f_*$,
\[
f_* \;:\; \mathcal{M} \longrightarrow \mathcal{N}\;,
\tag*{(4.1.2)}
\]
which therefore commutes with small $\varprojlim$\,. The datum of $f_*$ satisfying this property therefore determines $f^*$ up to unique isomorphism. But in general, it is not enough to give oneself such an $f_*$, to be able to assert that it corresponds to an $f^*$, i.e. that it admits a right adjoint\,; for this, one would have to know that the functors $\mathcal{M}\longrightarrow\Ens$ which commute with small $\varprojlim$ are representable, which will be the case if $\mathcal{M}$ admits a small \textit{co}generating subcategory, i.e. if $\mathcal{M}^\circ$ is also a pseudo-topos category. % [sic] qu'il admette un adjoint à droite -> adjoint à gauche, f^* is the left adjoint of f_*

\newpage

\printedpage{2}
% head: A. Grothendieck : Dérivateurs, Ch. XVIII, Édition des Universités de Montpellier II et Paris VII
% foot: Version 17 décembre 2025

It is no doubt judicious to proceed as in the case of topoi, and to call \textit{morphism of pseudo-topoi} $f:\mathcal{X}\longrightarrow\mathcal{Y}$ a quadruple
\[
(f^*,f_*,\alpha,\beta)
\tag*{(4.1.3)}
\]
\msp{page 53}
of two adjoint functors $(f^*,f_*)$, and adjunction morphisms. When $f^*$ is an equivalence, which amounts to $f_*$ being an equivalence, or $\alpha$ and $\beta$ isomorphisms, we will say that we have an \textit{equivalence of pseudo-topoi}, and for all practical purposes, we can then regard $\mathcal{X}$ and $\mathcal{Y}$ as essentially identical objects.

For two pseudo-topoi $\mathcal{X}$, $\mathcal{Y}$, I feel like introducing (as in the case of topoi) a product pseudo-topos
\[
\mathcal{X}\times\mathcal{Y}\;,
\tag*{(4.1.4)}
\]
which has the properties of a 2-product in the 2-category of pseudo-topoi (\footnote{It finally appears that what one constructs is \textit{not} a 2-product in the 2-category under consideration, so it is better to denote it by a sign such as $\mathcal{X}\boxtimes\mathcal{Y}$. But if $\mathcal{X}$, $\mathcal{Y}$ are two topoi, $\mathcal{X}\boxtimes\mathcal{Y}$ is a 2-product in the 2-category of topoi.}). I should specify that by definition
\[
\begin{aligned}
\uHom_{\pstop}(\mathcal{X},\mathcal{Y}) \;&\simeq\; \uHom_!(\mathcal{F}(\mathcal{Y}),\mathcal{F}(\mathcal{X}))^\circ\\
&\simeq\; \uHom^{!'}(\mathcal{F}(\mathcal{X}),\mathcal{F}(\mathcal{Y})),
\end{aligned}
\tag*{(4.1.5)}
\]
where the subscript ! (resp. the superscript !) indicates the full subcategory in the $\uHom$ category under consideration formed by the functors which commute with small $\varinjlim$ (resp. with small
\msp{page 54}
$\varprojlim$\,), and the prime in the $\uHom^{!'}$ in the last member indicates that we require moreover that the functor in question have a left adjoint.

So let
\[
\mathcal{M} \;=\; \mathcal{F}(\mathcal{X})\;, \qquad \mathcal{N} \;=\; \mathcal{F}(\mathcal{Y})\;,
\tag*{(4.1.6)}
\]
and let $\mathcal{Z}$ be a third pseudo-topos, described by
\[
\mathcal{P} \;=\; \mathcal{F}(\mathcal{Z})\;.
\tag*{(4.1.7)}
\]
We want to make explicit, up to equivalence, the product category $\uHom_{\pstop}(\mathcal{Z},\mathcal{X})\times\uHom_{\pstop}(\mathcal{Z},\mathcal{Y})$, i.e.
\[
\uHom_!(\mathcal{M},\mathcal{P})^\circ \;\times\; \uHom_!(\mathcal{N},\mathcal{Q})^\circ.
\tag*{(4.1.8)}
\] % [sic] \mathcal{Q} -> \mathcal{P}, \mathcal{Q} is not defined; \mathcal{Z} is described by \mathcal{P} (4.1.7)
For this I am going to introduce a category
\printedpage{3}
% head: A. Grothendieck : Dérivateurs, Ch. XVIII, Édition des Universités de Montpellier II et Paris VII
% foot: Version 17 décembre 2025
\msp{page 55}
\[
\mathcal{M}\boxtimes\mathcal{N}
\tag*{(4.1.9)}
\]
by the formula
\[
\mathcal{M}\boxtimes\mathcal{N} \;\simeq\; \uHom^{!!}(\mathcal{M}^\circ\times\mathcal{N}^\circ,\Ens)\;,
\tag*{(4.1.10)}
\]
where the sign $\uHom^{!!}$ denotes the full subcategory of the $\uHom$ category, formed by the bi-contrafunctors % doubt resolved: retain the diplomatic bi-contrafoncteurs reading; see corrections.csv, PDF leaf 5
\[
F \;:\; \mathcal{M}^\circ\times\mathcal{N}^\circ \longrightarrow \Ens\;,
\tag*{(4.1.10.1)}
\]
such that for every $x$ in $\mathcal{M}$, the functor
\[
y \longmapsto F(x,y) \;:\; \mathcal{N}^\circ \longrightarrow \Ens
\tag*{(4.1.10.2)}
\]
be in $\uHom^!(\mathcal{N}^\circ,\Ens)$, i.e. transforms small $\varinjlim$ into $\varprojlim$\,, and that, symmetrically, for every $y$ in $\mathcal{N}$, the functor
\[
x \longmapsto F(x,y) \;:\; \mathcal{M}^\circ \longrightarrow \Ens
\tag*{(4.1.10.3)}
\]
be in $\uHom^!(\mathcal{M}^\circ,\Ens)$. But this also means that these functors are \textit{representable}\,; so we also find
\[
\mathcal{M}\boxtimes\mathcal{N} \xleftarrow{\sim} \uHom^!(\mathcal{N}^\circ,\mathcal{M})
\tag*{(4.1.11)}
\]
and
\[
\mathcal{M}\boxtimes\mathcal{N} \xleftarrow{\sim} \uHom^!(\mathcal{M}^\circ,\mathcal{N})
\tag*{(4.1.12)}
\]
\msp{page 56}
using (in each of these formulas) on two sides the equivalences of categories
\[
\left\{
\begin{aligned}
&\mathcal{M} &&\xrightarrow{\sim} &&\uHom^!(\mathcal{M}^\circ,\Ens)\\
&\mathcal{N} &&\xrightarrow{\sim} &&\uHom^!(\mathcal{N}^\circ,\Ens)\;.
\end{aligned}
\right.
\tag*{(4.1.13)}
\]
We can also say that the category $\mathcal{M}\boxtimes\mathcal{N}$ is equivalent to that of \textit{pairs} of functors
\[
\left\{
\begin{aligned}
&\varphi &&: &&\mathcal{M}^\circ &&\longrightarrow &&\mathcal{N}\\
&\psi &&: &&\mathcal{N}^\circ &&\longrightarrow &&\mathcal{M}\;,
\end{aligned}
\right.
\tag*{(4.1.14)}
\]
endowed with two adjunction morphisms
\[
\left\{
\begin{aligned}
&\psi\varphi^\circ &&\xleftarrow{\alpha} &&\mathrm{id}_{\mathcal{M}}\\
&\varphi\psi^\circ &&\xleftarrow{\beta} &&\mathrm{id}_{\mathcal{N}}
\end{aligned}
\right.
\tag*{(4.1.15)}
\]
satisfying the two usual compatibilities expressing the adjunction of $\varphi^\circ$ and $\psi$ (or, equivalently, of $\psi^\circ$ and $\varphi$). The functor
\[
(\varphi,\psi) \longmapsto \varphi \qquad (\text{resp.}\;\; (\varphi,\psi) \longmapsto \psi)
\]
gives an equivalence of the category of these pairs with $\uHom^!(\mathcal{M}^\circ,\mathcal{N})$ (resp. with $\uHom^!(\mathcal{N}^\circ,\mathcal{M})$), so that we have a commutative square of canonical equivalences of functors:

% ---- from der-ch18_p006-010.tex
\printedpage{4}
% head: A. Grothendieck : Dérivateurs, Ch. XVIII, Édition des Universités de Montpellier II et Paris VII

\msp{page 57}

\[
\begin{tikzcd}[column sep=small]
 & \cM \tbt \cN \arrow[dl, "\sim" sloped] \arrow[dr, "\sim" sloped] & \\
\uHom^{!}(\cM^{\circ},\cN) \arrow[dr, "\sim" sloped] & & \uHom^{!}(\cN^{\circ},\cM) \arrow[dl, "\sim" sloped] \\
 & \underbrace{\cM \bt \cN}_{\defeq \uHom^{!!}(\cM^{\circ}\times\cN^{\circ},\Ens)} &
\end{tikzcd}
\tag*{(4.1.16)}
\]

\texttt{[where $\cM \tbt \cN$ is the category of pairs defined on the preceding page.]}

I say that $\cM \bt \cN$ \textit{is a pseudo-topos category}. It is clear, \texttt{[more]}over, that it admits small $\varprojlim$ , which are computed argument by argument in $\uHom(\cM^{\circ}\times\cN^{\circ},\Ens)$, or as one prefers in $\uHom^{!}(\cM^{\circ},\cN)$, or $\uHom^{!}(\cN^{\circ},\cM)$. How to construct small $\varinjlim$ ? There is no question of invoking an argument of cogenerators. One must moreover also find a small strictly generating subcategory. Visibly, there is something nontrivial to prove!

\section*{4.2\quad Fumblings on $\cM \bt \cN$}

Let us return to the general situation of a pseudo-topos $\cX$, described by the pseudo-topos category $\cM$. If $\cN$ is an arbitrary category, we define the category

\msp{page 58}
\[
\mathrm{Fais}(\cX,\cN) \qquad \text{or} \qquad \cF(\cX,\cN)
\]
of sheaves on $\cX$, with values in $\cN$, by
\[
\cF(\cX,\cN) \ \defeq\ \uHom^{!}(\cM^{\circ},\cN) \ .
\tag*{(4.2.1)}
\]
We thus have a fully faithful functor
\[
\cF(\cX,\cN) \ \defeq\ \uHom^{!}(\cM^{\circ},\cN) \ \linj\ \Hom^{!}(\cN^{\circ},\overline{\cM}) \ ,
\tag*{(4.2.2)}
\]
where $\overline{\cM}$ denotes the subcategory of $\cM^{\wedge}$ (\textbf{N.B.} $\cM^{\wedge}$ is not in general a $\mathcal{U}$-category never mind that) formed by the \textit{representable} functors $\cM^{\circ}\longrightarrow\Ens$ (which is a $\mathfrak{U}$-category, equivalent to $\cM$ by the canonical inclusion $\cM \linj \overline{\cM}$). So we have up to unique isomorphism a factorization of (4.2.2) into
\[
\cF(\cX,\cN) \ \overset{\text{f. faithful}}{\lhook\joinrel\relbar\joinrel\relbar\joinrel\longrightarrow}\ \uHom^{!}(\cN^{\circ},\cM) \ ,
\tag*{(4.2.3)}
\]
the essential image being formed by the functors
\[
\psi \ :\ \cN^{\circ} \longrightarrow \cM
\]
such that for every $x$ in $\cM$,
\[
y \longmapsto \Hom(y,\psi(x)) \ :\ \cN^{\circ} \longrightarrow \Ens
\]
be representable, i.e. by the functors which
% foot: Version 17 décembre 2025
\printedpage{5}
% head: A. Grothendieck : Dérivateurs, Ch. XVIII, Édition des Universités de Montpellier II et Paris VII
\msp{page 59}
admit a \textit{left adjoint}. When $\cN$ itself is a pseudo-topos category, this condition is automatic (it follows from the hypothesis that $\psi$ commutes with small $\varprojlim$ ). So in this case, (4.2.3) is an equivalence of categories.

Now let $S$ ($S$ as ``site'') be a \textit{small} subcategory of $\cM$, generating by strict epimorphisms. We know that this means that the composed functor
\[
\cM \xrightarrow{\ \sim\ } \uHom^{!}(\cM^{\circ},\Ens) \xrightarrow{\text{restriction to } S} S^{\wedge} \ (\ \defeq \uHom(S^{\circ},\Ens)\ )
\tag*{(4.2.4)}
\]
is fully faithful, and this implies more generally, that for any category $\cN$, the restriction functor
\[
\cF(\cX,\cN) \ \defeq\ \uHom^{!}(\cM^{\circ},\cN) \ \overset{\text{f. faithful}}{\lhook\joinrel\relbar\joinrel\relbar\joinrel\longrightarrow}\ \uHom(S^{\circ},\cN)
\tag*{(4.2.5)}
\]
is fully faithful. Indeed, the second member, via the fully faithful embedding
\[
\cN \ \overset{\text{f. faithful}}{\lhook\joinrel\relbar\joinrel\relbar\joinrel\longrightarrow}\ \uHom^{!}(\cN^{\circ},\Ens) \ \linj\ \cN^{\wedge} \ ,
\]
embeds into
\[
\uHom(S^{\circ},\uHom^{!}(\cN^{\circ},\Ens)) \ \simeq\ \Hom^{!}(\cN^{\circ},S^{\wedge}) \ ,
\]
and the full faithfulness of (4.2.5) is therefore equivalent to that of the composed functor
\[
\cF(\cX,\cN) \ \defeq\ \uHom^{!}(\cM^{\circ},\cN) \ \longrightarrow\ \uHom^{!}(\cN^{\circ},S^{\wedge}) \ \simeq\ \uHom(S^{\circ},\uHom^{!}(\cN^{\circ},\Ens)) \ .
\tag*{(4.2.6)}
\]

\msp{page 60}
Now we have commutativity in the square
\[
\begin{tikzcd}[column sep=large, row sep=large]
\cF(\cX,\cN) \defeq \uHom^{!}(\cM^{\circ},\cN) \arrow[r, "\text{\scriptsize(4.2.5)}"] \arrow[d, hook, "\text{\scriptsize(4.2.3)}"', "\text{\scriptsize fully faithful}"] & \uHom(S^{\circ},\cN) \arrow[d, hook, "\text{\scriptsize fully faithful}"] \\
\uHom^{!}(\cN^{\circ},\cM) \arrow[r, hook, "\text{\scriptsize(4.2.4)}", "\text{\scriptsize fully faithful}"'] & \underbrace{\uHom^{!}(\cN^{\circ},S^{\wedge})}_{\simeq\, \uHom(S^{\circ},\uHom^{!}(\cN^{\circ},\Ens))}
\end{tikzcd}
\tag*{(4.2.7)}
\]
(where the vertical arrows are equivalences, if $\cN$ is itself a pseudo-topos category). Let us note on the other hand that we have an equivalence
\[
\uHom(S^{\circ},\cN) \ \xleftarrow{\ \sim\ }\ \uHom^{!}(S^{\wedge\circ},\cN) \ \defeq\ \cF(\Top(S),\cN) \ ,
\tag*{(4.2.8)}
\]
where $\Top(S)$ is the pseudo-topos (in fact, a topos) defined by the pseudo-topos category $S^{\wedge}$. For this, let us note that we have
\[
\uHom^{!}(S^{\wedge\circ},\cN) \ \xrightarrow{\ \sim\ }\ \uHom_{!}(S^{\wedge},\cN^{\circ})^{\circ} \ ,
\]
and on the other hand the restriction functor to $S \linj S^{\wedge}$ gives
\[
\uHom_{!}(S^{\wedge},\cN^{\circ}) \ \xrightarrow{\ \sim\ }\ \uHom(S,\cN^{\circ}) \ ,
\]
whence
% foot: Version 17 décembre 2025